% !TeX root = ../Book5.tex
\begin{figure}[htbp]
\centering
\begin{tikzpicture}[x=0.85cm,y=0.85cm,every node/.style={font=\small}]
  % A_n: the dashed segment is a path, with no intermediate vertex when n=3.
  \node at (2,1.45) {\(A_n\quad(n\geq3)\)};
  \draw[thick] (0,0)--(1,0);
  \draw[thick,densely dashed] (1,0)--(4,0);
  \foreach \x in {0,1,4} \fill (\x,0) circle (1.6pt);
  \node at (2,-0.55) {\(n\) 个顶点};

  % D_4 is shown separately: all three arms have exactly one vertex.
  \node at (9.5,1.45) {\(D_4\)};
  \draw[thick] (8.5,0)--(10.5,0) (9.5,0)--(9.5,0.85);
  \foreach \x/\y in {8.5/0,10.5/0,9.5/0.85}
    \fill (\x,\y) circle (1.6pt);
  \fill[AlgebraBlue] (9.5,0) circle (2.2pt);
  \node at (9.5,-0.55) {臂长 \((1,1,1)\)};

  % D_n: the long arm has n-3 vertices; at n=5 its dashed segment is one edge.
  \node at (2,-1.25) {\(D_n\quad(n\geq5)\)};
  \draw[thick] (0,-2.7)--(2,-2.7) (1,-2.7)--(1,-1.85);
  \draw[thick,densely dashed] (2,-2.7)--(4,-2.7);
  \foreach \x/\y in {0/-2.7,2/-2.7,4/-2.7,1/-1.85}
    \fill (\x,\y) circle (1.6pt);
  \fill[AlgebraBlue] (1,-2.7) circle (2.2pt);
  \node at (2,-3.25) {臂长 \((1,1,n-3)\)};

  \node at (9.5,-1.25) {\(E_6\)};
  \draw[thick] (7.5,-2.7)--(11.5,-2.7) (9.5,-2.7)--(9.5,-1.85);
  \foreach \x in {7.5,8.5,10.5,11.5}
    \fill (\x,-2.7) circle (1.6pt);
  \fill (9.5,-1.85) circle (1.6pt);
  \fill[AlgebraBlue] (9.5,-2.7) circle (2.2pt);
  \node at (9.5,-3.25) {臂长 \((1,2,2)\)};

  \node at (2.5,-3.95) {\(E_7\)};
  \draw[thick] (0,-5.4)--(5,-5.4) (2,-5.4)--(2,-4.55);
  \foreach \x in {0,1,3,4,5} \fill (\x,-5.4) circle (1.6pt);
  \fill (2,-4.55) circle (1.6pt);
  \fill[AlgebraBlue] (2,-5.4) circle (2.2pt);
  \node at (2.5,-5.95) {臂长 \((1,2,3)\)};

  \node at (9.5,-3.95) {\(E_8\)};
  \draw[thick] (6.5,-5.4)--(12.5,-5.4) (8.5,-5.4)--(8.5,-4.55);
  \foreach \x in {6.5,7.5,9.5,10.5,11.5,12.5}
    \fill (\x,-5.4) circle (1.6pt);
  \fill (8.5,-4.55) circle (1.6pt);
  \fill[AlgebraBlue] (8.5,-5.4) circle (2.2pt);
  \node at (9.5,-5.95) {臂长 \((1,2,4)\)};
\end{tikzpicture}
\caption[ADE 图及臂长]{ADE 图及臂长。蓝色顶点为三叉点，臂长计中心之外的顶点数。
虚线表示链段，其中可以没有额外顶点：\(A_3\) 与 \(D_5\) 的虚线各是一条边。
\(A_1\) 是单顶点，\(A_2\) 是单边；\(D_4\) 的三条臂均长 \(1\)。}
\label{b5:fig:ade-diagrams}
\end{figure}
